\documentclass[UTF8,11pt]{ctexart}
\usepackage[a4paper,margin=24mm]{geometry}
\usepackage{amsmath,amssymb,amsthm,mathtools,mathrsfs}
\usepackage[all]{xy}
\usepackage{xurl,hyperref,enumitem}
\usepackage{tikz}
\usetikzlibrary{arrows.meta,cd,decorations.markings,calc,patterns}
\hypersetup{hidelinks,breaklinks=true}
\setlength{\emergencystretch}{3em}
\allowdisplaybreaks
\newtheorem*{theorem}{Theorem}
\newtheorem*{thm}{Theorem}
\newtheorem*{lemma}{Lemma}
\newtheorem*{proposition}{Proposition}
\newtheorem*{prop}{Proposition}
\newtheorem*{corollary}{Corollary}
\newtheorem*{cor}{Corollary}
\newtheorem*{claim}{Claim}
\theoremstyle{definition}
\newtheorem*{definition}{Definition}
\newtheorem*{defn}{Definition}
\newtheorem*{remark}{Remark}
\newtheorem*{example}{Example}
\newcommand{\R}{\mathbb R}
\newcommand{\C}{\mathbb C}
\newcommand{\Z}{\mathbb Z}
\newcommand{\Q}{\mathbb Q}
\newcommand{\N}{\mathbb N}
\newcommand{\F}{\mathbb F}
\newcommand{\D}{\mathbb D}
\newcommand{\bB}{\mathbb B}
\newcommand{\bC}{\mathbb C}
\newcommand{\bR}{\mathbb R}
\newcommand{\bZ}{\mathbb Z}
\newcommand{\bN}{\mathbb N}
\newcommand{\bQ}{\mathbb Q}
\newcommand{\bD}{\mathbb D}
\newcommand{\bF}{\mathbb F}
\newcommand{\CP}{\mathbb{CP}}
\newcommand{\RP}{\mathbb{RP}}
\newcommand{\sO}{\mathscr O}
\newcommand{\sI}{\mathscr I}
\newcommand{\sF}{\mathscr F}
\newcommand{\sG}{\mathscr G}
\newcommand{\sH}{\mathscr H}
\newcommand{\sR}{\mathscr R}
\newcommand{\sM}{\mathscr M}
\newcommand{\sV}{\mathscr V}
\newcommand{\sU}{\mathscr U}
\DeclareMathOperator{\Hom}{Hom}
\DeclareMathOperator{\Ext}{Ext}
\DeclareMathOperator{\Tor}{Tor}
\DeclareMathOperator{\Spec}{Spec}
\DeclareMathOperator{\Proj}{Proj}
\DeclareMathOperator{\supp}{supp}
\DeclareMathOperator{\im}{im}
\DeclareMathOperator{\id}{id}
\DeclareMathOperator{\Tr}{Tr}
\DeclareMathOperator{\tr}{tr}
\DeclareMathOperator{\rank}{rank}
\DeclareMathOperator{\ord}{ord}
\DeclareMathOperator{\dist}{dist}
\DeclareMathOperator{\End}{End}
\DeclareMathOperator{\Aut}{Aut}
\DeclareMathOperator{\Gal}{Gal}
\DeclareMathOperator{\vol}{vol}
\DeclareMathOperator{\Cl}{Cl}
\DeclareMathOperator{\coker}{coker}
\DeclareMathOperator{\codim}{codim}
\DeclareMathOperator{\divergence}{div}
\DeclareMathOperator{\Ric}{Ric}
\DeclareMathOperator{\grad}{grad}
\DeclareMathOperator{\Hess}{Hess}
\newcommand{\abs}[1]{\left\lvert#1\right\rvert}
\newcommand{\sabs}[1]{\lvert#1\rvert}
\newcommand{\babs}[1]{\bigl\lvert#1\bigr\rvert}
\newcommand{\Babs}[1]{\Bigl\lvert#1\Bigr\rvert}
\newcommand{\bbabs}[1]{\biggl\lvert#1\biggr\rvert}
\newcommand{\BBabs}[1]{\Biggl\lvert#1\Biggr\rvert}
\newcommand{\norm}[1]{\left\lVert#1\right\rVert}
\newcommand{\snorm}[1]{\lVert#1\rVert}
\newcommand{\bnorm}[1]{\bigl\lVert#1\bigr\rVert}
\newcommand{\Bnorm}[1]{\Bigl\lVert#1\Bigr\rVert}
\newcommand{\bbnorm}[1]{\biggl\lVert#1\biggr\rVert}
\newcommand{\BBnorm}[1]{\Biggl\lVert#1\Biggr\rVert}
\newcommand{\widebar}[1]{\overline{#1}}
\newcommand{\nicefrac}[2]{\frac{#1}{#2}}
\newcommand{\linnprod}[2]{\langle#1,#2\rangle}
\newcommand{\myindex}[1]{#1}
\newcommand{\glsadd}[1]{}
\newcommand{\avoidbreak}{}
\newcommand{\sourceref}[1]{\textnormal{[原文：\nolinkurl{#1}]}}
\newcommand{\thmref}[1]{\sourceref{#1}}
\newcommand{\lemmaref}[1]{\sourceref{#1}}
\newcommand{\propref}[1]{\sourceref{#1}}
\newcommand{\corref}[1]{\sourceref{#1}}
\newcommand{\defnref}[1]{\sourceref{#1}}
\newcommand{\exampleref}[1]{\sourceref{#1}}
\newcommand{\exerciseref}[1]{\sourceref{#1}}
\newcommand{\figureref}[1]{\sourceref{#1}}
\newcommand{\sectionref}[1]{\sourceref{#1}}
\newcommand{\appendixref}[1]{\sourceref{#1}}
\newcommand{\stacksref}[2]{\href{https://stacks.math.columbia.edu/tag/#1}{#2 [#1]}}


\newcommand{\E}{\mathbb E}
\newcommand{\Prob}{\mathbb P}
\newcommand{\Reff}{\mathcal R_{\mathrm{eff}}}
\newcommand{\eps}{\varepsilon}
\DeclareMathOperator{\diam}{diam}
\DeclareMathOperator{\Var}{Var}
\DeclareMathOperator{\Crit}{Crit}
\newcommand{\dd}{\,\mathrm d}


\begin{document}
\section*{077\quad 一致算术理论的Rosser不完备性}
\subsection*{来源与整理范围}
Open Logic Project，\emph{The Open Logic Text}，Incompleteness / incompleteness-provability，rosser-thm.tex（Git blob b4cf6667cb2611abc9d83f17c3f03891b5ab462a）；支撑 fixed-point-lemma.tex（aee58c2581b18e8757df96e8a3e1043339aeb0d2）。
\par\url{https://github.com/OpenLogicProject/OpenLogic/blob/master/content/incompleteness/incompleteness-provability/rosser-thm.tex}
\par\textbf{恢复方式（整理者说明）：}依据人类原文重新整理以补齐缺失题号；并非旧题证文件的逐字节恢复；2026-09-28。
\subsection*{作者命题与人类证明}

\paragraph{Notation (editor's expansion of the source macros).}
$\mathrm Q$ denotes Robinson arithmetic. For a formula $A$, $\ulcorner A\urcorner$ is the standard numeral of its G\"odel number; $\operatorname{diag}$ is the primitive recursive diagonal-substitution function. The formula $D_{\operatorname{diag}}(x,y)$ represents this function in $\mathrm Q$.
\begin{lemma}[Fixed-point lemma; the author's support proof]
Let $B(x)$ be any formula with one free variable $x$. Then there is a sentence $A$ such that $\mathrm Q\vdash A\leftrightarrow B(\ulcorner A\urcorner)$.
\end{lemma}
\begin{proof}
Given $B(x)$, let $E(x)$ be the formula $\exists y(D_{\operatorname{diag}}(x,y)\land B(y))$ and let $A$ be its diagonalization, i.e., the formula $E(\ulcorner E(x)\urcorner)$.
Since $D_{\operatorname{diag}}$ represents $\operatorname{diag}$, and $\operatorname{diag}(\ulcorner E(x)\urcorner)=\ulcorner A\urcorner$, $\mathrm Q$ can derive
\begin{align}
&D_{\operatorname{diag}}(\ulcorner E(x)\urcorner,\ulcorner A\urcorner),\label{repdiag1}\\
&\forall y(D_{\operatorname{diag}}(\ulcorner E(x)\urcorner,y)\to y=\ulcorner A\urcorner).\label{repdiag2}
\end{align}
Now we show that $\mathrm Q\vdash A\leftrightarrow B(\ulcorner A\urcorner)$. We argue informally, using just logic and facts derivable in $\mathrm Q$.
First, suppose $A$, i.e., $E(\ulcorner E(x)\urcorner)$. Going back to the definition of $E(x)$, we see that $E(\ulcorner E(x)\urcorner)$ just is
\[\exists y(D_{\operatorname{diag}}(\ulcorner E(x)\urcorner,y)\land B(y)).\]
Consider such a $y$. Since $D_{\operatorname{diag}}(\ulcorner E(x)\urcorner,y)$, by \eqref{repdiag2}, $y=\ulcorner A\urcorner$. So, from $B(y)$ we have $B(\ulcorner A\urcorner)$.
Now suppose $B(\ulcorner A\urcorner)$. By \eqref{repdiag1}, we have
\[D_{\operatorname{diag}}(\ulcorner E(x)\urcorner,\ulcorner A\urcorner)\land B(\ulcorner A\urcorner).\]
It follows that
\[\exists y(D_{\operatorname{diag}}(\ulcorner E(x)\urcorner,y)\land B(y)).\]
But that's just $E(\ulcorner E(x)\urcorner)$, i.e., $A$.
\end{proof}

\begin{theorem}[Rosser's theorem; the counted problem]
Let $T$ be any consistent, axiomatizable theory extending $\mathrm Q$. Then $T$ is not complete.
\end{theorem}
\begin{proof}
Recall that $\operatorname{Prov}_T(y)$ is defined as $\exists x\operatorname{Prf}_T(x,y)$, where $\operatorname{Prf}_T(x,y)$ represents the decidable relation which holds iff $x$ is the G\"odel number of a derivation of the sentence with G\"odel number $y$. The relation that holds between $x$ and $y$ if $x$ is the G\"odel number of a \emph{refutation} of the sentence with G\"odel number $y$ is also decidable. Let $\operatorname{not}(x)$ be the primitive recursive function which does the following: if $x$ is the code of a formula $A$, $\operatorname{not}(x)$ is a code of $\neg A$. Then $\operatorname{Refut}_T(x,y)$ holds iff $\operatorname{Prf}_T(x,\operatorname{not}(y))$. Let the corresponding arithmetic formula represent it. Then, if $T\vdash\neg A$ and $\delta$ is a corresponding derivation, $\mathrm Q\vdash\operatorname{Refut}_T(\ulcorner\delta\urcorner,\ulcorner A\urcorner)$. We define $\operatorname{RProv}_T(y)$ as
\[\exists x\bigl(\operatorname{Prf}_T(x,y)\land\forall z(z<x\to\neg\operatorname{Refut}_T(z,y))\bigr).\]
Roughly, $\operatorname{RProv}_T(y)$ says ``there is a proof of $y$ in $T$, and there is no shorter refutation of $y$.'' Assuming $T$ is consistent, $\operatorname{RProv}_T(y)$ is true of the same numbers as $\operatorname{Prov}_T(y)$; but from the point of view of \emph{provability} in $T$ (and we now know that there is a difference between truth and provability!) the two have different properties. If $T$ is \emph{in}consistent, then the two do \emph{not} hold of the same numbers!
By the fixed-point lemma, there is a formula $R_T$ such that
\begin{equation}\mathrm Q\vdash R_T\leftrightarrow\neg\operatorname{RProv}_T(\ulcorner R_T\urcorner).\label{RT}\end{equation}
In contrast to the proof of the first incompleteness theorem, here we claim that if $T$ is consistent, $T$ doesn't derive $R_T$, and $T$ also doesn't derive $\neg R_T$. (In other words, we don't need the assumption of $\omega$-consistency.)
First, let's show that $T\nvdash R_T$. Suppose it did, so there is a derivation of $R_T$ from $T$; let $n$ be its G\"odel number. Then $\mathrm Q\vdash\operatorname{Prf}_T(\bar n,\ulcorner R_T\urcorner)$, since $\operatorname{Prf}_T$ represents the proof relation in $\mathrm Q$. Also, for each $k<n$, $k$ is not the G\"odel number of a derivation of $\neg R_T$, since $T$ is consistent. So for each $k<n$, $\mathrm Q\vdash\neg\operatorname{Refut}_T(\bar k,\ulcorner R_T\urcorner)$. By the bounded-numeral lemma \textsf{lem:less-nsucc},
\[\mathrm Q\vdash\forall z(z<\bar n\to\neg\operatorname{Refut}_T(z,\ulcorner R_T\urcorner)).\]
Thus,
\[\mathrm Q\vdash\exists x\bigl(\operatorname{Prf}_T(x,\ulcorner R_T\urcorner)\land\forall z(z<x\to\neg\operatorname{Refut}_T(z,\ulcorner R_T\urcorner))\bigr),\]
but that's just $\operatorname{RProv}_T(\ulcorner R_T\urcorner)$. By \eqref{RT}, $\mathrm Q\vdash\neg R_T$. Since $T$ extends $\mathrm Q$, also $T\vdash\neg R_T$. We've assumed that $T\vdash R_T$, so $T$ would be inconsistent, contrary to the assumption of the theorem.
Now, let's show that $T\nvdash\neg R_T$. Again, suppose it did, and suppose $n$ is the G\"odel number of a derivation of $\neg R_T$. Then $\operatorname{Refut}_T(n,\ulcorner R_T\urcorner)$ holds, and since the corresponding formula represents it in $\mathrm Q$, $\mathrm Q\vdash\operatorname{Refut}_T(\bar n,\ulcorner R_T\urcorner)$. We'll again show that $T$ would then be inconsistent because it would also derive $R_T$. Since
\[\mathrm Q\vdash R_T\leftrightarrow\neg\operatorname{RProv}_T(\ulcorner R_T\urcorner),\]
and since $T$ extends $\mathrm Q$, it suffices to show that
\[\mathrm Q\vdash\neg\operatorname{RProv}_T(\ulcorner R_T\urcorner).\]
The sentence $\neg\operatorname{RProv}_T(\ulcorner R_T\urcorner)$, i.e.,
\[\neg\exists x\bigl(\operatorname{Prf}_T(x,\ulcorner R_T\urcorner)\land\forall z(z<x\to\neg\operatorname{Refut}_T(z,\ulcorner R_T\urcorner))\bigr),\]
is logically equivalent to
\[\forall x\bigl(\operatorname{Prf}_T(x,\ulcorner R_T\urcorner)\to\exists z(z<x\land\operatorname{Refut}_T(z,\ulcorner R_T\urcorner))\bigr).\]
We argue informally using logic, making use of facts about what $\mathrm Q$ derives. Suppose $x$ is arbitrary and $\operatorname{Prf}_T(x,\ulcorner R_T\urcorner)$. We already know that $T\nvdash R_T$, and so for every $k$, $\mathrm Q\vdash\neg\operatorname{Prf}_T(\bar k,\ulcorner R_T\urcorner)$. Thus, for every $k$ it follows that $x\ne\bar k$. In particular, we have (a) that $x\ne\bar n$. We also have $\neg(x=\bar0\lor x=\bar1\lor\cdots\lor x=\overline{n-1})$ and so by \textsf{lem:less-nsucc}, (b) $\neg(x<\bar n)$. By \textsf{lem:trichotomy}, $\bar n<x$. Since $\mathrm Q\vdash\operatorname{Refut}_T(\bar n,\ulcorner R_T\urcorner)$, we have $\bar n<x\land\operatorname{Refut}_T(\bar n,\ulcorner R_T\urcorner)$, and from that
\[\exists z(z<x\land\operatorname{Refut}_T(z,\ulcorner R_T\urcorner)).\]
Since $x$ was arbitrary we get, as required, that
\[\forall x\bigl(\operatorname{Prf}_T(x,\ulcorner R_T\urcorner)\to\exists z(z<x\land\operatorname{Refut}_T(z,\ulcorner R_T\urcorner))\bigr).\]
\end{proof}

\subsection*{整理说明与先修范围（非作者正文）}
由此前已读取的源文件正文恢复整理，保留Rosser构造及两个不可证方向，附入作者不动点引理的完整证明。开头记号段为编辑宏展开说明；源文件将元层关系和算术公式分别用Prf/OPrf、Refut/ORefut表示，此处以所在语境区分，未改变两个层次。删去原文括号中的shmovable用语旁论，不删推导。axiomatizable依原书指可有效公理化；$\bar n$是标准数码。标准先修为可计算关系在Robinson算术中的可表示性、固定标准数码的有限界枚举及三歧引理；不把Rosser构造或两向矛盾作为先修。难点是同时比较证明与反证编码，将仅有一致性转化为两个方向的不可证性。
\subsection*{可修改的简短标注（非作者正文）}
$c$：有效算术理论的完备性限制。

$a$：语法算术化；可表示关系；对角不动点；证明与反证编码比较；标准数码有限界。

\texttt{cross\_theory\_status = yes}。把关于证明存在与否的元数学命题编码为算术谓词，再以对角句和有限编码比较返回理论的不可判定句。位置：式RT及其后两个方向。

全部标注为非穷尽、可修改初稿。
\subsection*{来源许可与编辑归属}
Open Logic Project，\emph{The Open Logic Text}，CC BY 4.0；许可见仓库README：\url{https://github.com/OpenLogicProject/OpenLogic/blob/master/README.md}。正文来源宏展开为标准LaTeX；编辑题号、取材范围与标注另列。
ChatGPT 加入题号、中文来源说明、整理范围和初稿标注；编辑说明与人类证明分列。
\end{document}
